\documentclass[10pt]{article}
\usepackage[utf8]{inputenc}
\usepackage[vietnamese]{babel}
\usepackage{amsmath,amssymb,stmaryrd}
\usepackage{graphicx}
\usepackage{xcolor}
\usepackage{tcolorbox}
\tcbuselibrary{skins}
\usepackage{paracol}
\usepackage{fancyhdr}
\usepackage[letterpaper, top=20mm, bottom=22mm, left=18mm, right=18mm]{geometry}

% Màu sắc nhận diện chuẩn sách Stewart Calculus
\definecolor{stewartblue}{RGB}{0, 130, 185}
\definecolor{thmred}{RGB}{205, 50, 50}
\definecolor{graytext}{RGB}{80, 80, 80}

% Định nghĩa ký hiệu kết thúc ví dụ / lời giải
\newcommand{\exampleend}{\hfill{\color{stewartblue}\rule{2mm}{2mm}}}

% Thiết lập Header & Footer
\pagestyle{fancy}
\fancyhf{}
\renewcommand{\headrulewidth}{0pt}
\renewcommand{\footrulewidth}{0pt}
\fancyhead[L]{\small\textsf{\textbf{MỤC 2.3}\quad Tính các giới hạn bằng các quy tắc giới hạn}}
\fancyhead[R]{\small\textsf{\textbf{101}}}
\fancyfoot[C]{%
  \fontsize{6.5pt}{8pt}\selectfont\color{gray}%
  Copyright 2021 Cengage Learning. All Rights Reserved. May not be copied, scanned, or duplicated, in whole or in part. WCN 02-200-203\\
  Due to electronic rights, some third party content may be suppressed from the eBook and/or eChapter(s). Editorial review has deemed that any suppressed content does not materially affect the overall learning experience. Cengage Learning reserves the right to remove additional content at any time if subsequent rights restrictions require it.%
}

% Cấu hình hộp định lý
\newtcolorbox{stewartbox}{
  enhanced,
  colback=white,
  colframe=thmred,
  arc=0mm,
  boxrule=0.7pt,
  left=7pt,
  right=7pt,
  top=6pt,
  bottom=6pt
}

\begin{document}

% Thiết lập tỷ lệ 2 cột bất đối xứng: Cột lề trái ~28%, cột nội dung chính ~72%
\setcolumnwidth{48mm, 122mm}
\setlength{\columnsep}{10mm}

\begin{paracol}{2}

% ========================================================
% KHỐI 1: LỜI GIẢI TIẾP THEO & HÌNH 5
% ========================================================
{\small\noindent Trong Ví dụ 2.4.4 đã chỉ ra rằng $\lim_{x \to 0^+} \sqrt{x} = 0$.
\vspace{8pt}

\begin{center}
  \includegraphics[width=\linewidth]{images/p0136_fig1.png}
  \vspace{2pt}
  \centerline{\small\textsf{\textbf{HÌNH 5}}}
\end{center}
}

\switchcolumn

\noindent{\bfseries\color{stewartblue}LỜI GIẢI}\quad Vì $f(x) = \sqrt{x - 4}$ với $x > 4$, ta có
\[
\lim_{x \to 4^+} f(x) = \lim_{x \to 4^+} \sqrt{x - 4} = \sqrt{4 - 4} = 0
\]
Vì $f(x) = 8 - 2x$ với $x < 4$, ta có
\[
\lim_{x \to 4^-} f(x) = \lim_{x \to 4^-} (8 - 2x) = 8 - 2 \cdot 4 = 0
\]
Các giới hạn bên phải và bên trái bằng nhau. Do đó giới hạn tồn tại và
\[
\lim_{x \to 4} f(x) = 0
\]
Đồ thị của $f$ được thể hiện trong Hình 5.\exampleend

% ========================================================
% KHỐI 2: VÍ DỤ 10 & HÀM PHẦN NGUYÊN (HÌNH 6)
% ========================================================
\switchcolumn*

\vspace{4pt}
{\small\noindent Các ký hiệu khác của $\llbracket x \rrbracket$ là $[x]$ và $\lfloor x \rfloor$. Hàm phần nguyên đôi khi còn được gọi là \textit{hàm sàn} (floor function).
\vspace{8pt}

\begin{center}
  \includegraphics[width=\linewidth]{images/p0136_fig2.png}
  \vspace{2pt}
  \centerline{\small\textsf{\textbf{HÌNH 6}}}
  \centerline{\small Hàm phần nguyên}
\end{center}
}

\switchcolumn

\vspace{4pt}
\noindent{\bfseries\color{stewartblue}VÍ DỤ 10}\quad \textbf{Hàm phần nguyên} (greatest integer function) được định nghĩa bởi $\llbracket x \rrbracket = $ số nguyên lớn nhất nhỏ hơn hoặc bằng $x$. (Chẳng hạn, $\llbracket 4 \rrbracket = 4$, $\llbracket 4{,}8 \rrbracket = 4$, $\llbracket \pi \rrbracket = 3$, $\llbracket \sqrt{2} \rrbracket = 1$, $\llbracket -\frac{1}{2} \rrbracket = -1$.) Chứng minh rằng $\lim_{x \to 3} \llbracket x \rrbracket$ không tồn tại.
\medskip

\noindent{\bfseries\color{stewartblue}LỜI GIẢI}\quad Đồ thị của hàm phần nguyên được thể hiện trong Hình 6. Vì $\llbracket x \rrbracket = 3$ với $3 \leqslant x < 4$, ta có
\[
\lim_{x \to 3^+} \llbracket x \rrbracket = \lim_{x \to 3^+} 3 = 3
\]
Vì $\llbracket x \rrbracket = 2$ với $2 \leqslant x < 3$, ta có
\[
\lim_{x \to 3^-} \llbracket x \rrbracket = \lim_{x \to 3^-} 2 = 2
\]
Bởi vì các giới hạn một phía này không bằng nhau, $\lim_{x \to 3} \llbracket x \rrbracket$ không tồn tại theo Định lý 1.\exampleend

% ========================================================
% KHỐI 3: ĐỊNH LÝ KẸP & HÌNH 7
% ========================================================
\switchcolumn*

\vspace*{36mm}
\begin{center}
  \includegraphics[width=\linewidth]{images/p0136_fig3.png}
  \vspace{2pt}
  \centerline{\small\textsf{\textbf{HÌNH 7}}}
\end{center}

\switchcolumn

\vspace{6pt}
\noindent{\color{stewartblue}\rule{2.3mm}{2.3mm}}\hspace{6pt}{\large\textbf{Định lý kẹp}}\par\vspace{4pt}
Hai định lý sau đây mô tả mối liên hệ giữa các giới hạn của các hàm số khi giá trị của hàm này lớn hơn (hoặc bằng) giá trị của hàm kia. Các chứng minh của chúng có thể tìm thấy trong Phụ lục F.
\smallskip

% Hộp định lý 2
\begin{stewartbox}
  \noindent\tcbox[colback=white, colframe=thmred, size=small, arc=1pt, boxrule=0.7pt, on line]{\textbf{\color{thmred}2}}\hspace{5pt}\textbf{\color{thmred}Định lý}\hspace{8pt}Nếu $f(x) \leqslant g(x)$ khi $x$ ở gần $a$ (ngoại trừ có thể tại $a$) và các giới hạn của $f$ và $g$ đều tồn tại khi $x$ tiến dần đến $a$, thì
  \[
  \lim_{x \to a} f(x) \leqslant \lim_{x \to a} g(x)
  \]
\end{stewartbox}
\smallskip

% Hộp định lý 3: Định lý kẹp
\begin{stewartbox}
  \noindent\tcbox[colback=white, colframe=thmred, size=small, arc=1pt, boxrule=0.7pt, on line]{\textbf{\color{thmred}3}}\hspace{5pt}\textbf{\color{thmred}Định lý kẹp}\hspace{8pt}Nếu $f(x) \leqslant g(x) \leqslant h(x)$ khi $x$ ở gần $a$ (ngoại trừ có thể tại $a$) và
  \[
  \lim_{x \to a} f(x) = \lim_{x \to a} h(x) = L
  \]
  thì
  \[
  \lim_{x \to a} g(x) = L
  \]
\end{stewartbox}
\smallskip

Định lý kẹp (The Squeeze Theorem), đôi khi còn được gọi là \textit{Định lý kẹp bánh mì} (Sandwich Theorem) hay \textit{Định lý kẹp chặt} (Pinching Theorem), được minh họa trong Hình 7. Định lý nói rằng nếu $g(x)$ bị kẹp giữa $f(x)$ và $h(x)$ ở gần $a$, và nếu $f$ cùng $h$ có chung giới hạn $L$ tại $a$, thì $g$ buộc phải có cùng giới hạn $L$ tại $a$.

\end{paracol}

\end{document}